% TOP_PROBLEM: {"id": 3302, "title": "Universal highly arc transitive digraphs", "category_id": 3, "category": {"id": 3, "name": "graph_theory", "display_name": "Graph Theory", "description": "Problems involving graphs, networks, and their properties.", "slug": "graph-theory", "order_index": 3, "created_at": "2026-07-31T15:26:25.671Z"}, "status": "open"}
% FIRST_POSTED: null
% ENTRY_KIND: statement_only
% Imported from: batch03.tex; UnsolvedMath ID 3302
\documentclass[11pt]{article}
\usepackage[margin=25mm]{geometry}
\usepackage{fontspec}
\setmainfont{Latin Modern Roman}
\usepackage{amsmath,amssymb,mathtools}
\usepackage{unicode-math}
\setmathfont{Latin Modern Math}
\usepackage{xurl,hyperref}
\hypersetup{colorlinks=true,urlcolor=blue}
\setcounter{secnumdepth}{0}
\setlength{\parindent}{0pt}
\setlength{\parskip}{5pt}
\setlength{\emergencystretch}{3em}
\newcommand{\sref}[2]{\hyperlink{source:#1:#2}{[S#2]}}
\newcommand{\cataloguescope}[1]{\paragraph{Recorded target.}#1}
\begin{document}
% UnsolvedMath ID 3302; source catalogue code OPG-677
\setcounter{section}{3301}
\section{Universal highly arc transitive digraphs}\label{3302}
{\small\sffamily UnsolvedMath ID 3302 · Graph Theory}
\subsection{Definitions and mathematical statement}

\paragraph{Shared notation.}$\mathbb N=\{1,2,\ldots\}$, and $\mathbb Z,\mathbb Q,\mathbb R,\mathbb C$ denote the integers, rationals, reals and complex numbers. Any other conventions are specified locally.


An $s$-arc of a digraph is a sequence $(x_0,\ldots,x_s)$ such that $x_i\to x_{i+1}$ is an arc for $0\leq i<s$ and $x_{i-1}\ne x_{i+1}$ for $0<i<s$. Highly arc-transitive means that the automorphism group acts transitively on $s$-arcs for every $s\geq0$. Locally finite means finite indegree and outdegree at every vertex. A walk is alternating if, at each interior vertex, its two incident successive arcs either both enter or both leave that vertex. Reachability is universal if every pair of arcs lies on some common alternating walk. Let $P$ denote the directed graph with vertex set $\mathbb Z$ and arcs $i\to i+1$; a digraph homomorphism preserves arcs.
\paragraph{Statement recorded in the supplied source.}
Does there exist an infinite, locally finite highly arc-transitive digraph having at least one arc and universal reachability? The source additionally asks whether an infinite highly arc-transitive digraph having at least one arc can have no homomorphism onto $P$, including with local finiteness; its discussion reports affirmative constructions for this subsidiary question. The cited later primary paper answers the universal-reachability question affirmatively as well. \sref{3302}{2} \sref{3302}{3}
\subsection{Short English statement}
Can an infinite digraph with at least one arc have all finite directed-path symmetries and finite degrees while every two arcs communicate through an alternating walk? The accompanying question concerns the absence of a homomorphism onto the directed integer path.
\subsection{Sources}
\begin{description}
\item[\hypertarget{source:3302:1}{S1}] ULAM.AI and the UnsolvedMath Contributors, UnsolvedMath: A Curated Collection of Open Mathematics Problems, record 3302 (OPG-677). CC BY 4.0. \url{https://huggingface.co/datasets/ulamai/UnsolvedMath}.
\item[\hypertarget{source:3302:2}{S2}] Open Problem Garden, Universal highly arc transitive digraphs, node 677, original problem page. Original question and full discussion. \url{https://www.openproblemgarden.org/op/universal_highly_arc_transitive_digraphs}.
\item[\hypertarget{source:3302:3}{S3}] Matt DeVos, Bojan Mohar and Robert Šámal, Highly arc-transitive digraphs -- counterexamples and structure (2011), abstract. \url{https://arxiv.org/abs/1110.2945}.
\end{description}
\label{end:3302}
\end{document}
